\documentclass[12pt]{article}

\usepackage{amsmath}
\usepackage{bm}
\usepackage{caption}
\usepackage{cleveref}
\usepackage{graphicx}
\usepackage{rsfso}

\graphicspath{{graphs}}

\title{\textbf{Extension 2 HSC Assessment 3}\\ {\Large Applied Calculus}}
\author{Zac Hart}

\begin{document}

\maketitle
\newpage

\begin{enumerate}
    % 1.
    \item \begin{enumerate}
              % 1.a.
              \item \textbf{For which values of $\bm{r}$ does the function $\bm{y = e^{r x}}$ satisfy the differential equation $\bm{y'' + y' - y = 0$} ?} \smallskip

                    Substituting for $y$ gives
                    \begin{align*}
                        r^2 e^{r x} + r e^{r x} - e^{r x} & = 0  \\
                        (r^2 + r - 1) e^{r x}             & = 0  \\
                        r^2 + r - 1                       & = 0.
                    \end{align*}
                    This is a quadratic in terms of $r$, so we can use the quadratic formula to get
                    \begin{align*}
                        r & = \frac{-1 \pm \sqrt{1 - (-4)}}{2} \\
                          & =\frac{-1 \pm \sqrt{5}}{2}.
                    \end{align*}

                    % 1.b.
              \item \textbf{If $\bm{r_1}$ and $\bm{r_2}$ are the values of $\bm{r}$ that you found in part (a), show that every member of the family of functions $\bm{y = a e^{r_1 x} + b e^{r_2 x}}$ is also a solution.} \smallskip

                    Again substituting for $y$,
                    \begin{align*}
                        \text{LHS} & = (a r_1^2 e^{r_1 x} + b r_2^2 e^{r_2 x}) + (a r_1 e^{r_1 x} + b r_2 e^{r_2 x}) - (a e^{r_1 x} + b e^{r_2 x}) \\
                                   & = (a r_1^2 e^{r_1 x} + a r_1 e^{r_1 x} - a e^{r_1 x}) + (b r_2^2 e^{r_2 x} + b r_2 e^{r_2 x} - b e^{r_2 x})   \\
                                   & = 0 + 0                                                                                                       \\
                                   & = \text{RHS},
                    \end{align*}
                    hence $y = a e^{r_1 x} + b e^{r_2 x}$ is a solution to the differential equation.
          \end{enumerate}

    \item \textbf{Which of the following functions are solutions of the differential equation $\bm{y'' + y = \sin x}$ ?}
          \begin{enumerate}
              \item $\bm{y = \sin x}$

                    \begin{align*}
                        \text{LHS} & = -\sin x + \sin x \\
                                   & = 0                \\
                                   & \neq \text{RHS},
                    \end{align*}
                    therefore $y = \sin x$ is not a solution.

              \item $\bm{y = \cos x}$

                    \begin{align*}
                        \text{LHS} & = -\cos x + \cos x \\
                                   & = 0                \\
                                   & \neq \text{RHS},
                    \end{align*}
                    therefore $y = \cos x$ is not a solution.

              \item $\bm{y = \frac{1}{2}x \sin x}$

                    \begin{align*}
                        \text{LHS} & = \frac{d^2}{dx^2} \tfrac{1}{2}x \sin x                              \\
                                   & = \tfrac{1}{2}\frac{d}{dx}[\sin x + x \cos x] + \tfrac{1}{2}x \sin x \\
                                   & = \tfrac{1}{2}(\cos x + \cos x - x \sin x) + \tfrac{1}{2}x \sin x    \\
                                   & = \cos x - \tfrac{1}{2}x \sin x + \tfrac{1}{2}x \sin x               \\
                                   & = \cos x                                                             \\
                                   & \neq \text{RHS},
                    \end{align*}
                    therefore $y = \sin x$ is not a solution.
              \item $\bm{y = -\frac{1}{2}x \cos x}$

                    \begin{align*}
                        \text{LHS} & = \frac{d^2}{dx^2}[-\tfrac{1}{2}x \cos x]                             \\
                                   & = -\tfrac{1}{2}\frac{d}{dx}[\cos x - x \sin x] + \tfrac{1}{2}x \cos x \\
                                   & = -\tfrac{1}{2}(-\sin x - \sin x - x \cos x) - \tfrac{1}{2}x \cos x   \\
                                   & = \sin x + \tfrac{1}{2}x \cos x - \tfrac{1}{2}x \cos x                \\
                                   & = \sin x                                                              \\
                                   & = \text{RHS},
                    \end{align*}
                    therefore $y = \sin x$ is a solution.
          \end{enumerate}

          % 3.
    \item \begin{enumerate}
              % 3.a.
              \item \textbf{Prove that if $\bm{f}$ is a continuous function, then}
                    \begin{equation*}
                        \bm{\int_0^a f(x) dx = \int_0^a f(a - x) dx.}
                    \end{equation*} \smallskip

                    Setting $u = a - x$, the RHS becomes
                    \begin{align*}
                        \int_a^0 -f(u) du
                         & = \int_0^a f(u) du \\
                         & = \text{RHS}.
                    \end{align*}

                    % 3.b.
              \item \textbf{Use the result in part (a) to show that}
                    \begin{equation*}
                        \bm{\int_0^\frac{\pi}{2} \frac{\sin^n x}{\sin^n x + \cos^n x} dx = \frac{\pi}{4}.}
                    \end{equation*} \smallskip

                    Let $I = \int_0^\frac{\pi}{2} \frac{\sin^n x}{\sin^n x + \cos^n x} dx$. \medskip

                    Using part (a) gives
                    \begin{align*}
                        I   & = \int_0^\frac{\pi}{2} \frac{\sin^n(\frac{\pi}{2} - x)}{\sin^n(\frac{\pi}{2} - x) + \cos^n(\frac{\pi}{2} - x)} dx \\
                            & = \int_0^\frac{\pi}{2} \frac{\cos^n x}{\cos^n x + \sin^n x} dx,
                        \intertext{hence}
                        2 I & = I + \int_0^\frac{\pi}{2} \frac{\cos^n x}{\cos^n x + \sin^n x} dx                                                \\
                            & = \int_0^\frac{\pi}{2} \frac{\sin^n x + \cos^n x}{\sin^n x + \cos^n x} dx                                         \\
                            & = \int_0^\frac{\pi}{2} dx                                                                                         \\
                            & = \frac{\pi}{2}.
                    \end{align*}
                    Thus $I = \frac{\pi}{4}$.
          \end{enumerate}

          % 4.
    \item \begin{enumerate}
              % 4.a.
              \item \textbf{By rotating the line $\bm{y = \frac{r}{h} x}$ about the $\bm{x}$-axis, show that the volume of the cone, height $\bm{h}$ is given by the formula $\bm{V = \frac{1}{3} \pi r^2 h}$ where $\bm{r}$ is the radius of the circular base and $\bm{h}$ is the vertical height of the cone.} \smallskip

                    \begin{figure}[ht]
                        \centering
                        \includegraphics[width=0.5\textwidth]{graph-1}
                        \caption{The surface of revolution.}
                        \label{fig:1}
                    \end{figure}

                    The surface of revolution (\cref{fig:1}) is a cone with radius $r$ when $x = h$. We are therefore interested in the volume under this surface from $x = 0$ to $x = h$.

                    From the formula
                    \begin{equation*}
                        V = \pi \int_a^b y^2 dx \\
                    \end{equation*}
                    we have
                    \begin{align*}
                        V & = \pi \int_0^h (\frac{r}{h} x)^2 dx                 \\
                          & = \frac{\pi r^2}{h^2} \int_0^h x^2 dx               \\
                          & = \frac{\pi r^2}{h^2} \cdot \frac{x^3}{3} \Big|_0^h \\
                          & = \frac{\pi r^2}{h^2} \frac{h^3}{3}                 \\
                          & = \tfrac{1}{3} \pi r^2 h.
                    \end{align*}

                    % 4.b.
              \item \textbf{If the semi-vertical angle of the apex of the cone is $\bm{\theta}$, show that the volume of the cone becomes $\bm{V = \frac{1}{3} \pi h^3 \tan^2{\theta}}$.} \smallskip

                    The angle $\theta$ is labeled in \cref{fig:2}. Using trigonometry gives the relation $r = h \tan{\theta}$, hence by substituting into the equation of $V$,
                    \begin{align*}
                        V & = \tfrac{1}{3} \pi r^2 h                \\
                          & = \tfrac{1}{3} \pi (h \tan{\theta})^2 h \\
                          & = \tfrac{1}{3} \pi h^3 \tan^2{\theta}.
                    \end{align*}

                    \begin{figure}[ht]
                        \centering
                        \includegraphics[width=0.5\textwidth]{graph-2}
                        \caption{A triangular cross-section of the cone.}
                        \label{fig:2}
                    \end{figure}
          \end{enumerate}

          % 5.
    \item \textbf{We know how to find the volume of a solid of revolution obtained by rotating a region about a horizontal or vertical line. However, what if we rotate about a slanted line, that is, a line that is neither horizontal nor vertical? In this question you are asked to discover formulas for the volume of a solid of revolution when the axis of rotation is a slanted line. Let $\bm{y = f(x)}$ be the "upper" boundary of the region to be rotated for $\bm{[p, q]}$ and let $\bm{y = mx + b}$ be the line about which the rotation is to occur.}
          \begin{enumerate}
              % 5.a.
              \item \textbf{Show that the area of the region $\mathcal{R}$ to be rotated is given by}
                    \begin{equation*}
                        \bm{\frac{1}{1 + m^2} \int^p_q (f(x) - m x - b)(1 + m f'(x)) dx.}
                    \end{equation*}

                    \begin{figure}[ht]
                        \centering
                        \includegraphics[width=0.7\textwidth]{graph-3.pdf}
                        \caption{The region $\mathcal{R}$.}
                        \label{fig:3}
                    \end{figure}

                    Integration is usually performed along the $x$ or $y$ axis, however we want to integrate along the line $y = m x + b$. This means that the idea of integration must be fundamentally re-evaluated.

                    The area under the continuous function $f(x)$ in an interval is approximated by summing all of the rectangles with width $\Delta x_i$ and height $f(x_i)$ in said interval. As the number of rectangles approaches infinity, the error of the sum approaches zero until the sum exactly equals the area under the curve.

                    Likewise, the area bounded by $f(x)$ and the line $y = m x + b$ is approximated by summing all the rectangles with width $\Delta u_i$ and height $f(u_i)$, then taking the limit as $\Delta u_i$ approaches zero (see \cref{fig:3}).

                    We begin by expressing $\Delta u_i$ in terms of $\Delta x_i$.

                    \begin{figure}[ht]
                        \centering
                        \includegraphics[width=0.9\textwidth]{graph-4.pdf}
                        \caption{The tangent line at a point on the curve.}
                        \label{fig:4}
                    \end{figure}

                    \begin{figure}[ht]
                        \centering
                        \includegraphics[width=0.9\textwidth]{graph-5.pdf}
                        \caption{A close-up view of \cref{fig:4}.}
                        \label{fig:5}
                    \end{figure}
          \end{enumerate}

\end{enumerate}
\end{document}
